\chapter{Chapter 1}

motivations...

\section{Problems of the Standard Model}

The most important theory of physics, with predictions that have been verified to an extremely high degree of accuracy, is the Standard Model of particle physics. It describes the electromagnetic, weak, and strong interactions between elementary particles.

Anomalous magnetic moment of the electron, $a_e$, is one of the most precisely measured quantities in physics. We look at a magnetic dipole moment \(\bs{\mu} = g\frac{q}{2m}\mathbf{s}\) where at tree level, \(g = 2\). The anomalous magnetic moment is defined as \(a_e = \frac{g-2}{2}\) and can be computed for electrons as well as for muons etc. The experimental value of \(a_e\) is known to an accuracy of 0.28 parts per billion:
\[
    a_e = 0.001159
\]
while in theory, the value is computed to be
\[
    a_e = 0.001159652181643(764)
\]
around ten digits of accuracy.

There is also evidence for physics beyond the Standard Model, as proven in LHC experiments. The Standard Model does not explain the existence of dark matter, which is known to exist from astrophysical observations. It also does not explain the matter-antimatter asymmetry in the universe, as well as the hierarchy problem, which is the question of why the Higgs boson mass is so much lighter than the Planck scale.

There are problems of different natures in the Standard Model:
\begin{itemize}
    \item \textbf{Incompleteness}: The Standard Model does not include gravity, which is described by General Relativity. It also does not account for dark matter and dark energy, which make up most of the universe's mass-energy content.
    \item \textbf{Technical problems}: something that is contained in the Standard Model but is not explained by it.
    \item \textbf{Naturalness issues}: there are some things that can be explained by the Standard Model, but they are not natural and if you perturb the parameters of the Standard Model, they would not be stable.
    \item \textbf{Why questions}: why are there 4 dimensions, why are there 3 generations of fermions, why is the Higgs mass so small, etc.
\end{itemize}
We will see what supersymmetry can do to address some of these problems in the next sections.

\subsection{Incompleteness}

Somethings are incomplete theoretically, others are incomplete experimentally.

\subsubsection{Quantum Gravity}

gravity is not renormalizable perturbatively, and we do not have a quantum theory of gravity. The Standard Model does not include gravity, which is described by General Relativity. we weould need infinitely many counterterms to renormalize gravity, and we do not have a quantum theory of gravity. We do not know how to quantize gravity, and we do not know how to unify it with the other forces. The good part is that gravity is very weak, so we can ignore it at low energies. The bad part is that we do not know what happens at high energies, and we do not know how to unify gravity with the other forces. Newtons constant is the coupling constant of gravity, and it has dimensions of inverse mass squared, while it should be dimensionless in order for the theory to be renormalizable. The Planck scale is the energy scale at which gravity becomes strong, and it is around \(10^{19}\) GeV (which is orders of magnitude higher than the energies we can probe in experiments). At that energy scale gravity becomes strongly coupled, and we need a quantum theory of gravity to describe it. Thus the standard model is good as an effective field theory at energies \(E \ll 10^{18} GeV\), such that it is possible to speak about a cutoff \(\Lambda \leq M_p\) and the theory is valid up to that scale. The Standard Model is not valid at energies above the Planck scale, and we need a more fundamental theory to describe physics at those energies. String theory is one candidate for a quantum theory of gravity, but it is not yet experimentally verified.

Consder now the \textbf{Einstein-Hilbert action} for gravity:
\[
    S = \int \dots
\]
and if one adds the standard model action to it, one gets vertices of interaction between gravitons and standard model particles. The graviton is the quantum of the gravitational field, and it is a massless spin-2 particle. The vertices of interaction between gravitons and standard model particles are suppressed by powers of the Planck scale, which means that they are very weak at low energies. However, at high energies, these interactions become strong, and we need a quantum theory of gravity to describe them. The \textbf{Euler Lagrange equations} for a metric \(g_{\mu\nu}(x)\) are given by
\[
    R_{\mu\nu} - \frac{1}{2}g_{\mu\nu}R + \Lambda g_{\mu\nu} = \frac{T_{\mu\nu}}{M_p^2} = -\frac{2}{M_p^2 \sqrt{-g}}\frac{\delta (\sqrt{-g} \mathcal{L}_{\text{SM}})}{\delta g^{\mu\nu}}
\]
whch can now be expanded around flat space \(g_{\mu\nu} = \eta_{\mu\nu} + 2\frac{h_{\mu\nu}(x)}{M_p}\) where \(h_{\mu\nu}\) is the graviton field. We can get an interaction lagrangian of the form
\[
    \mathcal{L}_{\text{int}} = \frac{1}{M_p}h_{\mu\nu}T^{\mu\nu},
\]
which is why we say that the graviton couples to the energy-momentum tensor of the Standard Model and the vertices are suppressed by the Planck scale. This comes from
\[
    \sqrt{-g} \mathcal{L}_{\text{SM}} = \sqrt{-\eta} \frac{\delta (\sqrt{-g} \mathcal{L}_{\text{SM}})}{\delta g^{\mu\nu}} \delta g^{\mu\nu}\dots
\]
\begin{example}[Scalar field coupled to gravity]
    Consider a scalar field \(\phi\) with energy-momentum tensor \(T_{\mu\nu} = \partial_\mu \phi^* \partial_\nu \phi + \partial_\nu \phi^* \partial_\mu \phi - g_{\mu\nu}(\partial^\alpha \phi^* \partial_\alpha \phi - m^2 \phi^* \phi)\). The interaction lagrangian is then ...

        [...]

    Which corresponds to the vertex of interaction between the graviton and the scalar field, which is suppressed by the Planck scale.
\end{example}

What can be a solution to go beyond the Standard Model and include gravity? One idea is supergravity, which is a theory that combines supersymmetry and general relativity. In supergravity, some dangerous divergences cancel out due to supersymmetry, and the theory is renormalizable at one loop. However, supergravity is not a complete theory of quantum gravity, and it is not yet experimentally verified. Another idea is string theory, which is a theory that describes particles as one-dimensional objects called strings. In string theory, gravity is naturally included, and the theory is finite and consistent at all energy scales. However, string theory is also not yet experimentally verified. But looking at the theory, from its foundations there is no underlying reason why it should be still renormalizable at higher loops, and it is not clear if it is a complete theory of quantum gravity.

Looking at string theory, the finite dimension of the string is the reason why it is finite, and it is not clear if it is a complete theory of quantum gravity. The string has a finite size, and it cannot probe distances smaller than its size. This means that the theory is not sensitive to the ultraviolet divergences that plague quantum field theories.

Other attempts ...

\subsubsection{No Neutrino Masses}

We measured neutrino oscillations at Gran Sasso and SuperKamiokande, which implies that neutrinos have mass. However, in the Standard Model, neutrinos are massless. This is a problem because it means that the Standard Model is incomplete and cannot explain the observed phenomena of neutrino oscillations. Not a big issue, \(m_\nu \sim 10^{-3} \) eV. The answer is to add right-handed neutrinos, which are singlets under the Standard Model gauge group, and give them a Majorana mass term. If they are Dirac neutrinos, then we can add a Yukawa coupling to the Higgs field, which gives them a mass after electroweak symmetry breaking. If they are Majorana neutrinos, then we can add a Majorana mass term, which violates lepton number conservation. This is known as the seesaw mechanism, and it can explain why neutrino masses are so small compared to other fermion masses... hyerarchy problems in coupling constants, etc.

\subsubsection{Dark Matter}

[...]

Dark matter has to be 25 percent of the energy density of the universe...

One solution can be to study new particles should be
\begin{enumerate}
    \item chargeless
    \item feebly interacting, like neutrinos, or interact through a new force
    \item stable, and we have to be carefull with the mass of the particle, because if it is too light, it would be relativistic and would not form structures in the universe. If it is too heavy, it would decay too quickly.
    \item cold as in non-relativistic, because if it is relativistic, it would not form structures in the universe. If it is too heavy, it would decay too quickly.
\end{enumerate}

There are two main classes of models for dark matter:
\begin{itemize}
    \item \(m_{DM} \) is large and only weakly interacting: the neutralino in supersymmetry is a good candidate for this class of models. It is a stable particle that interacts only through the weak force, and it can have a mass in the range of 100 GeV to 1 TeV (cold since it is very heavy). It has spin 1/2, and it is a Majorana fermion. It is the lightest supersymmetric particle (LSP) in many supersymmetric models. it could be the supersimmetric partner of the photon, the Z boson, or the Higgs boson (spin 1 mediators). LSP furthermore are stable. Thermal quantum field theory (thermal equilibrium at the time of the early universe) predicts that the relic density of dark matter is inversely proportional to the annihilation cross section, which is of the order of the weak interaction. When the universe cools down, the dark matter particles freeze out and their number density becomes constant. Standard model particles decreases in numbers since they can annihilate, but they also freeze out as the universe expands.

              [...] WIMP miracle: the relic density of dark matter is of the order of the observed value if the annihilation cross section is of the order of the weak interaction. This is a coincidence, and it is known as the WIMP miracle. The WIMP miracle is a strong motivation for supersymmetry, because it predicts a stable particle with a mass in the range of 100 GeV to 1 TeV that interacts only through the weak force.

    \item completely different scenario: \(m_{DM} \) is small and ultra weak interaction: QCD axions are a good candidate for this class of models. CP symmetry is preserved in strong interactions, but it is violated in weak interactions. In QCD there is a term in the Lagrangian that violates CP symmetry, but experiments show that this term is very small, which is known as the strong CP problem. The axion is a hypothetical particle that was proposed to solve the strong CP problem. It is a pseudoscalar particle that weakly couples and has a very small mass. How can it be cold? It is produced non-thermally in the early universe (out of the standard model thermal equilibrium), and we can imagine a potential which is zero for \(E>\Lambda_{QCD}\), but for \(E<\Lambda_{QCD}\) it develops a potential with a minimum at \(\theta = 0\). The axion field is initially displaced from the minimum, and it starts to oscillate around the minimum, but then this classical field oscillation behaves like a cold dark matter fluid.... \(\phi_0, m_{\phi} \) ...
\end{itemize}